% ============================================================
% kbs_util.tex — KBS journal: Downstream utility (Section 6).
% Section label: sec:util. RQ1 (plan): representation-dependence via
% field ablation; completeness-dependence via the coverage sweep;
% robustness. Numbers macro-driven; E1 scale-up marked \TODO.
% Bibliography: kbs_ke_references.bib + paper/references.bib.
% ============================================================

\section{Downstream Utility of Acquired Knowledge}
\label{sec:util}

Acquisition validity (Section~\ref{sec:valid}) answers whether the acquired
knowledge is correct; utility answers whether a knowledge-based system that
consumes it performs better. The consuming system is the metadata-aware
retrieval index (Section~\ref{sec:method:system}); utility is measured as
retrieval quality, nDCG@10, over the discovery question pool with bootstrap
confidence intervals and paired permutation tests \cite{ndcg}. We report
\emph{per representation} (which knowledge field carries the effect) and
\emph{per completeness} (how the effect depends on how much knowledge the
base already has), on the full in-corpus question pool
\TODO{E1: scale to the 1,300-book corpus}.

\subsection{Retrieval gain of acquired knowledge}

Indexing machine-acquired subject knowledge into records that had none
raises topical retrieval substantially: on the conference-scale corpus,
the metadata-aware configuration reaches topical nDCG@10 of
\TopicalMeta{} against \TopicalHybrid{} for the best metadata-blind
configuration (paired permutation test \TopicalP{}; dense, the strongest
single metadata-blind mode, reaches \TopicalDense{})
\TODO{E1: re-verify at scale}. Overall retrieval rises from
\OverallHybrid{} to \OverallMeta{}. The gain is class-specific: on
known-item questions the metadata-aware configuration is nominally ahead
(\KnownMeta{} vs.\ \KnownHybrid{}) but not significantly (\KnownP{}),
and on bibliographic-fact questions it is significantly \emph{worse}
(\BibMeta{} vs.\ \BibHybrid{}, \BibP{}), a regression that reproduces on
the larger build \TODO{E1: report the extended-corpus replication}. The
utility of acquired knowledge is therefore concentrated in topical
discovery; we claim nothing for the other classes and recommend routing
bibliographic-fact queries around the record path.

\subsection{Representation dependence: which knowledge field carries the effect}

The central knowledge-engineering result is \emph{which} representation the
consuming system uses. We rebuild the record index eight times, varying only
which knowledge fields enter the record text (authors are always present as
the identity baseline), with retrieval, generation, and the chunk index
fixed. Adding subject headings to title is worth \AblSubjGain{} topical
nDCG@10 (95\% CI \AblSubjCI{}, \AblSubjP{}); adding classification is worth
\AblDdcGain{} (\AblDdcCI{}, \AblDdcP{})---a confidence interval that does
not reach a tenth of the subject effect and whose upper end is zero. The
full index is not better than title-plus-headings on topical questions
(\AblFullTopical{} vs.\ \AblTitleSubjTopical{}); the classification field
is, at best, inert for topical discovery. The two fields also serve
disjoint access modes: subject headings alone reach \AblSubjOnlyTopical{}
on topical questions but only \AblSubjOnlyKnown{} on known-item, while
title alone inverts that pattern (\AblTitleOnlyKnown{} known-item,
\AblTitleOnlyTopical{} topical). Authors alone reach \AblAuthorsTopical{}
on topical, confirming the topical signal is the acquired knowledge, not
corpus structure.

Classification has one reliable contribution, and it is an artifact of the
question class: on bibliographic-fact questions---which ask what class a
record is filed under---indexing classification is worth
\AblDdcBibGain{} (\AblDdcBibP{}). That is the system answering questions
\emph{about} the knowledge, not discovery \emph{by} it. We therefore index
title, author, and subject headings in the main configuration and exclude
classification---a decision this ablation licenses rather than assumes.

\subsection{Completeness dependence: how much knowledge the base already has}

Machine-acquired knowledge should matter most where the knowledge base has
gaps. We measure this directly by varying how many gold-bearing records keep
their headings and gap-filling every record left without one
(Section~\ref{sec:method:benchmark}), then reading topical retrieval
against coverage (conference scale; the sweep is the same construction the
journal corpus extends \TODO{E1}):

\begin{center}
\begin{tabular}{lccc}
\toprule
Knowledge-base coverage & no acquisition & gap-filled & gain \\
\midrule
0\% & \SpBaseZero{} & \SpEnrZero{} & +\SpGainZero{} \\
25\% & \SpBaseTwentyFive{} & \SpEnrTwentyFive{} & +\SpGainTwentyFive{} \\
50\% & \SpBaseFifty{} & \SpEnrFifty{} & +\SpGainFifty{} \\
75\% & \SpBaseSeventyFive{} & \SpEnrSeventyFive{} & +\SpGainSeventyFive{} \\
100\% & \SpBaseHundred{} & \SpEnrHundred{} & +\SpGainHundred{} \\
\bottomrule
\end{tabular}
\end{center}

The gain falls monotonically from +\SpGainZero{} at zero coverage to
+\SpGainHundred{} at full coverage (both columns differ at
\SpP{}), but never reaches zero: even a fully cataloged collection has
gaps worth filling. The best knowledge base in this study is neither human
nor machine alone---gap-filling at full coverage
(\SpEnrHundred{}) exceeds the professionally cataloged records on their
own (\SpBaseHundred{}) and exceeds an all-LLM catalog
\TODO{E1: re-verify loop conditions at scale}. And the benefit concentrates
exactly where it is needed: records that never received professional
headings are invisible to topical search until someone writes one, and it
does not have to be a cataloger.

\subsection{What the acquired knowledge is not: robustness}

Two checks confirm the mechanism is knowledge, not retrieval noise. First,
corrupting subject headings degrades discovery monotonically: topical
nDCG@10 falls from \NoiseZeroTopical{} (no corruption) to
\NoiseHundredTopical{} (100\% corruption, plausible-wrong or missing
headings), and overall from \NoiseZeroOverall{} to \NoiseHundredOverall{}.
Known-item drifts slightly \emph{up} under corruption (\NoiseZeroKnown{}
$\rightarrow$ \NoiseHundredKnown{}) and bibliographic-fact stays flat,
so the degradation is carried entirely by the subject-knowledge path---the
same path the field ablation isolates. At 100\% corruption, topical
retrieval falls \emph{below} the title-only baseline
(\NoiseHundredTopical{} vs.\ \AblTitleOnlyTopical{}): wrong knowledge is
worse than no knowledge, the expected ordering, and a cross-check that the
two experiments agree. Second, we implemented the natural alternative
retrieval design---fusing full-text chunk rankings into the record ranking
as a four-way reciprocal-rank fusion---and it is consistently worse for
discovery (\TopicalMeta{} $\rightarrow$ \FourTopical{} topical;
\OverallMeta{} $\rightarrow$ \FourOverall{} overall): mixing full-text
evidence into the record ranking dilutes the subject-knowledge signal that
topical discovery depends on.

\subsection{Scoring validity: single-gold bias and graded relevance}

Every topical number above is scored against the single record each question
was generated from, which understates systems that retrieve a relevant
non-seed record. We bound this with pooled graded relevance: for each
topical question the top records from all configurations are pooled and
judged (0/1/2) without knowledge of the seed
\TODO{E2: extend the judged pool at journal scale}. The pool is genuinely
multi-relevant---mean \PoolMeanRel{} fully-relevant records per query, with
\PoolMultiRel{} of \PoolNQueries{} questions having more than one---and the
effect on system ordering is small, with one exception. The gap-filling
margin shrinks under graded scoring: on the full topical pool,
single-gold scoring puts gap-filling +\PoolSingleEnrGain{} ahead of the
gold-only catalog, and pooled graded judging puts it +\PoolEnrGain{} ahead.
The direction is unchanged---gap-filling still wins---but the honest effect
size is smaller than the single-gold headline, and we report both. The
judgments are LLM-derived and are calibrated against human raters: two
raters graded a stratified sample of 120 pooled pairs, with human--human
agreement ($\kappa=0.755$, 3-way) and human--LLM agreement
($\kappa=0.64$--$0.80$) indicating the judge agrees with a human rater
about as well as human raters agree with each other
\TODO{E2: extend calibration to the full journal pool}.

\subsection{Summary}

Downstream utility has three measured properties. It is
\emph{representation-dependent}: subject headings carry the entire topical
gain (\AblSubjGain{}, \AblSubjP{}) while classification contributes none
(\AblDdcGain{}, CI upper end zero), so a knowledge-acquisition budget should
target the thesaurus representation. It is \emph{completeness-dependent}:
gain is monotone decreasing in existing coverage and never reaches zero, so
acquisition is a gap-filling channel whose marginal value a manager can read
off the coverage curve. And it is \emph{robust to the scoring choice} in
direction if not in magnitude (+\PoolSingleEnrGain{} single-gold,
+\PoolEnrGain{} graded), with the judge human-calibrated.
